%
%  chapter5-discussion.tex
%  ISELthesis
%
%  Discussion and Critical Analysis of the OmniFlow + QuickFaaS integration.
%
\chapter{Discussion and Critical Analysis}
\label{cha:discussion}

The previous chapter described \emph{what} was built. This chapter steps back to assess
\emph{how well} it was built, \emph{how} the design reached its current shape, what it would
mean to operate it inside a regulated retail bank, and \emph{where} it should go next. The
assessment is deliberately candid: every claim, favourable or otherwise, is anchored to a
concrete class or benchmark so that the reader can verify it against the source tree.

\section{Critical Assessment}
\label{sec:critical-assessment}

\subsection{Design and Extensibility}

The integration's central design decision --- abstracting function deployment behind the
\texttt{InternalFunctionDeployer} strategy interface --- is sound and pays off. The interface
exposes a single operation, \texttt{deployOrUpdate(functionName, deploymentDescriptorPath)},
and has exactly three implementations: \texttt{NoopInternalFunctionDeployer} (a safe default
that raises rather than silently doing nothing), \texttt{QuickFaasDeployer} (GCP), and
\texttt{AwsLambdaDeployer} (AWS). The two providers are kept structurally symmetric: each
provider deployer (\texttt{GoogleCloudDeployer}, \texttt{AmazonCloudDeployer}) owns a
provider-specific resolver (\texttt{WorkflowInternalFunctionResolver},
\texttt{AwsInternalFunctionResolver}) and both share the \texttt{FunctionRegistryStore}. Adding
a third platform therefore means adding one resolver and one deployer, without editing the
existing ones --- a genuine open/closed extension point rather than a claimed one.

The DSL boundary is equally disciplined. The mutual exclusivity of internal and external calls
is enforced \emph{twice}: once optimistically at build time in \texttt{CallContextBuilder}
(the builder refuses \texttt{host()}/\texttt{path()} once \texttt{internalFunction(...)} is set,
and vice versa) and again defensively at resolution time (both resolvers reject a
\texttt{CallContext} that carries an internal reference together with a non-blank host or path).
The resolution logic itself is expressed as a clean recursive tree walk that descends into
branches, iterations and parallel blocks, so an internal call nested arbitrarily deep is still
intercepted. Diagnostics are uniformly fail-fast and actionable --- an unresolved function, an
ambiguous regional match, or a stale registry entry each abort deployment with a message that
names the offending function and the corrective action.

\paragraph{Weakness: design residue left in the active path.}
The most visible design flaw is archaeological rather than functional. An earlier registry
model based on two synthetic keys per function (\texttt{"<functionRef>.host"} and
\texttt{"<functionRef>.path"}), embodied by \texttt{FunctionEndpointKeys} and
\texttt{WorkflowInternalCallEndpointResolver}, was superseded by the single
\texttt{serviceName}+\texttt{url} metadata model (\texttt{FunctionInvocationMetadata}) that the
two real resolvers consume. The superseded classes were never removed:
\texttt{WorkflowInternalCallEndpointResolver} is still \emph{imported} by
\texttt{GoogleCloudDeployer} yet never instantiated in the deployment path, and both it and
\texttt{FunctionEndpointKeys} survive only through their unit tests. This is dead code on the
critical path's doorstep. It is harmless at runtime but misleading to a maintainer, and --- as
noted in the review of Chapter~\ref{cha:impl} --- it is the source of a documentation drift, in
which the ``missing \texttt{functionRef}'' error semantics describe the abandoned model rather
than the shipped one.

\subsection{Performance and Scalability}

The performance story is the strongest part of the contribution, precisely because it follows
the full loop of \emph{identify $\rightarrow$ fix $\rightarrow$ quantify}. Benchmark P3 first
established that resolving an internal call costs roughly $7.5\,\mu$s and scales
linearly with the number of calls, because each call re-reads and re-parses the whole registry
file from disk. P6 isolated the second dimension and characterised the cost as growing in direct
proportion to both the number of calls and the number of registry entries (with $N$ calls over an
$R$-entry registry): about $180\,\mu$s of fixed I/O and parsing per read plus roughly
$0.56\,\mu$s per registered entry. P7 then applied the obvious remedy --- read the
registry once per workflow into an in-memory snapshot and resolve against it --- eliminating that
call-times-registry-size product in favour of a cost that no longer multiplies the two factors
together (one read proportional to the registry size, plus a cheap constant-time lookup per call),
delivering up to a \(\sim\!200\times\) speedup in the
worst corner ($N{=}200$, $R{=}1000$: \(\sim\!105\)\,ms down to \(\sim\!0.54\)\,ms). The
optimisation is real and lives in \texttt{WorkflowInternalCallEndpointResolver.resolve} and in
\texttt{FunctionRegistryStore.resolveUrlIn}.

\begin{table}[htbp]
  \centering
  \caption{P3 --- resolution cost of internal vs.\ external calls, $R{=}1$.}
  \label{tab:eval-p3}
  \begin{tabular}{rrr}
    \hline
    \textbf{Calls ($N$)} & \textbf{Internal (\(\mu\)s)} & \textbf{External (\(\mu\)s)} \\
    \hline
    1   & 7.5    & 0.02 \\
    10  & 74.6   & 0.14 \\
    50  & 368.0  & 0.65 \\
    100 & 747.2  & 1.10 \\
    200 & 1511.4 & 2.34 \\
    \hline
  \end{tabular}
\end{table}

\begin{table}[htbp]
  \centering
  \caption{P6 --- resolution cost per call as a function of registry size $R$, independent of $N$.}
  \label{tab:eval-p6}
  \begin{tabular}{rr}
    \hline
    \textbf{Registry size ($R$)} & \textbf{Cost per call (\(\mu\)s)} \\
    \hline
    1    & 181.0 \\
    10   & 185.1 \\
    50   & 207.0 \\
    200  & 279.2 \\
    1000 & 741.7 \\
    \hline
  \end{tabular}
\end{table}

\begin{table}[htbp]
  \centering
  \caption{P7 --- total resolution time (\(\mu\)s), naive (per-call read) vs.\ optimized (single read).}
  \label{tab:eval-p7}
  \begin{tabular}{lrrrr}
    \hline
     & $N{=}1$ & $N{=}10$ & $N{=}50$ & $N{=}200$ \\
    \hline
    naive, $R{=}1$        & 4.6   & 47.0   & 233.6   & 920.2    \\
    optimized, $R{=}1$    & 4.7   & 5.0    & 5.0     & 5.7      \\
    naive, $R{=}1000$     & 543.7 & 5480.8 & 26698.5 & 104889.0 \\
    optimized, $R{=}1000$ & 557.4 & 535.7  & 551.8   & 536.6    \\
    \hline
  \end{tabular}
\end{table}

This snapshot-based optimisation is not merely a property of the isolated
\texttt{FunctionRegistryStore} primitive: P8 exercises the resolution of a real \texttt{Workflow}
object across a grid of $1$ to $50$ distinct functions and $1$ to $200$ calls, and the same
pattern holds at the whole-workflow level --- the naive path stays dominated by the number of
calls and largely indifferent to how many distinct functions are involved (roughly $200\,\mu$s
per call at low $N$), while the optimised path stays close to flat, with the measured speedup
climbing from about $1\times$ at $N{=}1$ to about $110\times$ at $N{=}200$. P9 then isolates the
registry-size effect on the same real workflow, fixing $10$ distinct functions and $50$ calls
while the registry grows from $10$ to $1000$ entries: the naive cost climbs with the registry
size while the optimised cost stays essentially flat, so the speedup itself grows, from roughly
$41\times$ at $R{=}10$ to roughly $58\times$ at $R{=}1000$.

\begin{table}[htbp]
  \centering
  \caption{P8 --- naive resolution time (\(\mu\)s) of a real workflow, $F$ distinct functions
  $\times$ $N$ calls, $R{=}F$.}
  \label{tab:eval-p8-naive}
  \begin{tabular}{rrrrrr}
    \hline
    $F\backslash N$ & 1 & 5 & 10 & 50 & 200 \\
    \hline
    1  & 216 & 1\,118 & 2\,103 & 9\,544  & 38\,394 \\
    2  & 192 & 956    & 1\,902 & 9\,555  & 37\,950 \\
    5  & 189 & 995    & 1\,911 & 9\,663  & 38\,468 \\
    10 & 195 & 1\,132 & 2\,390 & 10\,544 & 39\,340 \\
    20 & 199 & 997    & 2\,096 & 10\,017 & 41\,205 \\
    50 & 219 & 1\,133 & 2\,194 & 10\,857 & 43\,360 \\
    \hline
  \end{tabular}
\end{table}

\begin{table}[htbp]
  \centering
  \caption{P8 --- optimized (single-read) resolution time (\(\mu\)s) of a real workflow, same
  grid as Table~\ref{tab:eval-p8-naive}.}
  \label{tab:eval-p8-optimized}
  \begin{tabular}{rrrrrr}
    \hline
    $F\backslash N$ & 1 & 5 & 10 & 50 & 200 \\
    \hline
    1  & 191 & 195 & 197 & 228 & 337 \\
    2  & 191 & 193 & 200 & 230 & 341 \\
    5  & 191 & 194 & 209 & 228 & 349 \\
    10 & 228 & 197 & 202 & 241 & 362 \\
    20 & 199 & 202 & 215 & 237 & 349 \\
    50 & 224 & 221 & 223 & 251 & 367 \\
    \hline
  \end{tabular}
\end{table}

\begin{table}[htbp]
  \centering
  \caption{P9 --- resolution time (\(\mu\)s) vs.\ registry size $R$, $F{=}10$/$N{=}50$ fixed.}
  \label{tab:eval-p9}
  \begin{tabular}{rrrr}
    \hline
    \textbf{$R$} & \textbf{naive} & \textbf{optimized} & \textbf{speedup} \\
    \hline
    10   & 9\,829  & 237 & 41.4$\times$ \\
    20   & 11\,952 & 308 & 38.8$\times$ \\
    50   & 10\,899 & 268 & 40.7$\times$ \\
    100  & 12\,118 & 275 & 44.1$\times$ \\
    200  & 14\,964 & 324 & 46.2$\times$ \\
    1000 & 43\,245 & 749 & 57.8$\times$ \\
    \hline
  \end{tabular}
\end{table}

The structural benchmarks reinforce the result: P18 and P19 show that nesting depth and parallel
width add no measurable cost beyond the number of leaf calls they contain, confirming that
resolution scales with the \emph{count} of internal calls, not the shape of the workflow tree.

\begin{table}[htbp]
  \centering
  \caption{P18 --- resolution time (\(\mu\)s) vs.\ nesting depth, 20 internal calls fixed,
  $R{=}F{=}10$.}
  \label{tab:eval-p18}
  \begin{tabular}{rr}
    \hline
    \textbf{Depth} & \textbf{Time (\(\mu\)s)} \\
    \hline
    0 & 207.3 \\
    1 & 203.2 \\
    2 & 200.7 \\
    3 & 200.4 \\
    4 & 204.1 \\
    5 & 199.6 \\
    \hline
  \end{tabular}
\end{table}

\begin{table}[htbp]
  \centering
  \caption{P19 --- resolution time (\(\mu\)s) vs.\ \texttt{Parallel} branch width, $N{=}5\times$
  width, $R{=}F{=}10$.}
  \label{tab:eval-p19}
  \begin{tabular}{rrr}
    \hline
    \textbf{Branches} & \textbf{$N$} & \textbf{Time (\(\mu\)s)} \\
    \hline
    1   & 5   & 188.3 \\
    2   & 10  & 195.0 \\
    5   & 25  & 204.5 \\
    10  & 50  & 220.8 \\
    20  & 100 & 276.2 \\
    50  & 250 & 404.9 \\
    100 & 500 & 583.7 \\
    \hline
  \end{tabular}
\end{table}

Two further benchmarks refine the resolution story in ways worth stating precisely. P14 measures
the fallback path used for regional disambiguation, where an exact-key lookup fails and the
resolver instead scans the already-loaded in-memory map for a \texttt{"region/functionRef"}
suffix match: this reintroduces a cost proportional to the registry size, but over data already
held in memory, so it stays cheap in absolute terms even as the ratio to an exact match grows to
about $1.9\times$ at a registry of $1000$ entries. P15 and P16 confirm that once a workflow mixes
internal and external calls, resolution cost tracks only the number of \emph{internal} calls, not
the workflow's total call count: a workflow with zero internal calls stays flat as the number of
external calls grows to $200$, because external calls never touch the registry, and the
registry-size effect measured in P9 is unchanged, within the same order of magnitude, once a
fraction of the calls are external.

\begin{table}[htbp]
  \centering
  \caption{P14 --- exact vs.\ suffix match (\(\mu\)s), $F{=}10$/$N{=}50$ fixed.}
  \label{tab:eval-p14}
  \begin{tabular}{rrrr}
    \hline
    \textbf{$R$} & \textbf{exact} & \textbf{suffix} & \textbf{ratio} \\
    \hline
    10   & 219 & 228    & 1.04$\times$ \\
    20   & 224 & 241    & 1.07$\times$ \\
    50   & 250 & 278    & 1.11$\times$ \\
    100  & 259 & 333    & 1.29$\times$ \\
    200  & 309 & 442    & 1.43$\times$ \\
    1000 & 768 & 1\,456 & 1.90$\times$ \\
    \hline
  \end{tabular}
\end{table}

\begin{table}[htbp]
  \centering
  \caption{P15 --- resolution time (\(\mu\)s) of a mixed workflow, $R{=}F{=}10$ fixed; $I$
  internal vs.\ $E$ external calls.}
  \label{tab:eval-p15}
  \begin{tabular}{rrrrrrr}
    \hline
    $I\backslash E$ & 0 & 1 & 5 & 10 & 50 & 200 \\
    \hline
    0   & 181 & 190 & 182 & 187 & 183 & 183 \\
    1   & 193 & 184 & 183 & 187 & 183 & 186 \\
    5   & 186 & 186 & 211 & 188 & 187 & 187 \\
    10  & 189 & 205 & 195 & 203 & 189 & 255 \\
    50  & 217 & 217 & 217 & 230 & 219 & 322 \\
    200 & 326 & 321 & 333 & 323 & 315 & 393 \\
    \hline
  \end{tabular}
\end{table}

\begin{table}[htbp]
  \centering
  \caption{P16 --- resolution time (\(\mu\)s) vs.\ registry size $R$, mixed workflow with
  $I{=}40$/$E{=}10$/$F{=}10$ fixed.}
  \label{tab:eval-p16}
  \begin{tabular}{rr}
    \hline
    \textbf{$R$} & \textbf{Time (\(\mu\)s)} \\
    \hline
    10   & 346 \\
    20   & 221 \\
    50   & 247 \\
    100  & 309 \\
    200  & 306 \\
    1000 & 776 \\
    \hline
  \end{tabular}
\end{table}

On the sizing axis, S1 measures the overhead that QuickFaaS's provider-agnostic layer adds to a
Lambda bundle at roughly $0.9$\,KB --- negligible for any real function, and in line
with the original QuickFaaS evaluation for GCP and Azure.

\begin{table}[htbp]
  \centering
  \caption{S1 --- Lambda bundle size, AWS agnostic layer vs.\ native implementation.}
  \label{tab:eval-s1}
  \begin{tabular}{lr}
    \hline
    \textbf{Bundle} & \textbf{Size} \\
    \hline
    AWS agnostic (QuickFaaS hook + \texttt{AwsHttpTemplate} + \texttt{aws-lambda-java-core}) & 14\,552 bytes \\
    AWS native (single \texttt{RequestHandler}, same logic) & 13\,647 bytes \\
    \textbf{Delta (agnostic-layer overhead)} & \textbf{905 bytes} \\
    \hline
  \end{tabular}
\end{table}

\paragraph{Weakness: the optimisation never reached production code.}
The single-read optimisation was demonstrated on the standalone store primitive, but the
resolvers that actually run during deployment --- \texttt{AwsInternalFunctionResolver} and
\texttt{WorkflowInternalFunctionResolver} --- were never converted. Benchmarks P10 and P11
measure exactly these classes, one per provider, and confirm they still pay the full cost,
growing in proportion to
both the number of calls and the registry size, on every deployment, because each call goes
through \texttt{FunctionRegistryStore.tryResolveEntry}, which re-reads the whole file. In other
words, the headline \(\sim\!200\times\) speedup is a property
of the benchmark harness, not of the shipping resolvers. This is acknowledged in the project's
own engineering notes, which flags it as ``a known optimization opportunity, not yet applied''.
It is a modest change --- read once in \texttt{resolve}, thread the snapshot through --- and its
absence is the clearest example of the prototype's reach exceeding its polish.

\begin{table}[htbp]
  \centering
  \caption{P10 --- AWS provider: resolution time (\(\mu\)s) of the real resolver
  (\texttt{AwsInternalFunctionResolver}), $R{=}F$.}
  \label{tab:eval-p10}
  \begin{tabular}{rrrrrr}
    \hline
    $F\backslash N$ & 1 & 5 & 10 & 50 & 200 \\
    \hline
    1  & 356 & 1\,684 & 3\,338 & 17\,806 & 71\,825 \\
    2  & 369 & 1\,744 & 3\,413 & 17\,205 & 83\,141 \\
    5  & 349 & 1\,730 & 3\,419 & 16\,961 & 67\,808 \\
    10 & 355 & 1\,738 & 3\,421 & 17\,509 & 69\,700 \\
    20 & 352 & 1\,779 & 3\,488 & 17\,455 & 69\,627 \\
    50 & 383 & 1\,847 & 3\,732 & 22\,396 & 73\,603 \\
    \hline
  \end{tabular}
\end{table}

\begin{table}[htbp]
  \centering
  \caption{P11 --- GCP provider: resolution time (\(\mu\)s) of the real resolver
  (\texttt{WorkflowInternalFunctionResolver}), $R{=}F$.}
  \label{tab:eval-p11}
  \begin{tabular}{rrrrrr}
    \hline
    $F\backslash N$ & 1 & 5 & 10 & 50 & 200 \\
    \hline
    1  & 398 & 2\,205 & 3\,837 & 19\,908 & 76\,812 \\
    2  & 400 & 1\,952 & 3\,959 & 20\,418 & 79\,856 \\
    5  & 459 & 2\,068 & 4\,001 & 24\,333 & 78\,215 \\
    10 & 395 & 1\,971 & 4\,436 & 19\,744 & 78\,525 \\
    20 & 402 & 2\,045 & 3\,955 & 20\,143 & 78\,037 \\
    50 & 428 & 2\,118 & 4\,285 & 20\,811 & 83\,269 \\
    \hline
  \end{tabular}
\end{table}

\paragraph{Weakness: write path scales poorly.}
Registration is worse than resolution. Every \texttt{put()} re-reads and rewrites the entire
registry file, so writing $K$ functions into a registry that already holds $R_0$ entries takes
longer than linear time in $K$ (P13): each successive write pays for a registry that has grown by
one entry since the write before it, so the cost compounds as $K$ grows, on top of the floor set
by $R_0$; the miss-then-write pattern the real resolvers use before
any cloud call adds a further 24--54\% on top (P17). For the small registries this work targets
this is irrelevant --- the absolute numbers stay in the hundreds of microseconds to low
hundreds of milliseconds, far below the seconds spent on the actual cloud deployment --- but it
is a genuinely non-linear write path hiding behind a convenient file abstraction.

\begin{table}[htbp]
  \centering
  \caption{P13 --- total write time (\(\mu\)s) for $K$ successive \texttt{put} calls, initial
  registry size $R_0$.}
  \label{tab:eval-p13}
  \begin{tabular}{rrrrrr}
    \hline
    $K\backslash R_0$ & 0 & 10 & 50 & 200 & 1000 \\
    \hline
    1   & 2\,673   & 2\,860   & 3\,235   & 4\,139   & 2\,761   \\
    5   & 12\,887  & 13\,149  & 14\,955  & 19\,710  & 13\,874  \\
    10  & 27\,691  & 26\,986  & 29\,912  & 39\,164  & 24\,029  \\
    50  & 154\,601 & 144\,606 & 161\,098 & 206\,385 & 123\,322 \\
    100 & 329\,757 & 304\,187 & 330\,100 & 419\,531 & 237\,910 \\
    \hline
  \end{tabular}
\end{table}

\emph{Note on Table~\ref{tab:eval-p13}:} the $R_0{=}1000$ column does not follow the otherwise
monotonic growth in $R_0$ at $K{=}100$; this single data point is attributed to noise on the
shared, non-dedicated benchmarking machine, not to a code effect.

\begin{table}[htbp]
  \centering
  \caption{P17 --- total time (\(\mu\)s) for $K$ successive miss+\texttt{put} pairs, initial
  registry size $R_0$.}
  \label{tab:eval-p17}
  \begin{tabular}{rrrrrr}
    \hline
    $K\backslash R_0$ & 0 & 10 & 50 & 200 & 1000 \\
    \hline
    1   & 3\,309   & 3\,545   & 3\,959   & 5\,122   & 4\,243   \\
    5   & 16\,699  & 18\,053  & 18\,795  & 23\,928  & 17\,403  \\
    10  & 35\,107  & 34\,358  & 37\,190  & 47\,631  & 47\,817  \\
    50  & 214\,238 & 185\,078 & 196\,112 & 250\,589 & 176\,340 \\
    100 & 410\,389 & 382\,140 & 403\,488 & 511\,878 & 346\,068 \\
    \hline
  \end{tabular}
\end{table}

\subsection{Robustness, Security and Operability}

Several properties that are acceptable in a proof of concept become liabilities under
production scrutiny.

\begin{itemize}
    \item \textbf{The registry is an unguarded flat file.} It has no in-memory cache, no
    integrity check (no signature, no checksum), and no cross-process locking --- only
    \texttt{@Synchronized} methods that serialise access \emph{within} a single JVM. Two
    concurrent deployments writing the same registry are a last-writer-wins race. More
    importantly for a security-sensitive setting, a tampered \texttt{url} or Lambda ARN in the
    file would be trusted and wired into the workflow without question; the registry is
    authoritative but unauthenticated.
    \item \textbf{IAM side effects are broad and automatic.} \texttt{AwsLambdaIamHelper} will,
    on its own initiative, create a shared \texttt{OmniFlowLambdaExecutionRole} attached to the
    managed \texttt{AWSLambdaBasicExecutionRole} policy, rewrite a role's trust policy when it
    finds \texttt{lambda.amazonaws.com} missing, and grant \texttt{states.amazonaws.com}
    permission to invoke the function. Each step is followed by a fixed \texttt{Thread.sleep}
    (10--30\,s) to absorb IAM eventual consistency. This is exactly the ergonomics a demo wants
    and exactly what a bank's least-privilege and change-control regime forbids: an automated
    tool mutating shared IAM roles outside the review pipeline.
    \item \textbf{The QuickFaaS boundary is a subprocess.} Deployment shells out to the
    QuickFaaS JAR via \texttt{QuickFaasProcessInvoker}, coupling two build systems (the Maven
    \texttt{deployment} module and the Gradle, Kotlin~1.6.20 QuickFaaS module) across a
    process boundary and a hard-coded JAR path. It works, but it is a brittle seam: failures
    surface as parsed subprocess output (the role-propagation retry, for instance, keys off the
    substring \texttt{"cannot be assumed by Lambda"} in the child's error text).
    \item \textbf{Credentials are environment-only.} Both providers authenticate through
    \texttt{EnvironmentVariableCredentialsProvider} (AWS) or ambient application-default
    credentials (GCP). There is no assume-role chain and no per-tenant isolation, so the tool
    runs with whatever ambient identity launched it.
\end{itemize}

\subsection{Testability}

The test suite is substantial (38 test classes on the \texttt{deployment} side) and the testing
philosophy is explicit and defensible: exercise only local logic --- rendering, resolution,
registry I/O --- and never call or mock the cloud SDKs. The consequence, however, is that the
riskiest code is the least covered. The classes that talk to AWS and GCP
(\texttt{AwsLambdaDeployer}, \texttt{AwsLambdaIamHelper}, \texttt{AwsRequests}, the Cloud
Functions HTTP calls) have no automated coverage at all, partly by choice and partly by
constraint: under the JDK~21 runtime \texttt{System.setSecurityManager} is unavailable, so code
paths that call \texttt{exitProcess} cannot be exercised in-process, and several GCP client
classes eagerly resolve application-default credentials merely on construction. The result is a
suite that thoroughly validates the pure core and leaves the integration edges --- where
production incidents actually happen --- untested.

\section{Evolution of the Approach}
\label{sec:evolution}

The source tree records its own history clearly enough that the design's evolution can be read
directly from it.

\paragraph{From two manual steps to one definition.}
The starting point was two separate, manually coordinated activities: deploy the functions with
QuickFaaS, then hand-wire their concrete endpoints into an OmniFlow workflow. The unification's
premise was to collapse these into a single workflow definition from which internal functions
are auto-deployed and their endpoints resolved automatically. The DSL change that anchors this
is the \texttt{internalFunction(name, deploymentDescriptorPath)} construct added to the
\texttt{call} step.

\paragraph{From literal endpoints to logical indirection.}
The first conceptual move was to replace literal \texttt{host}/\texttt{path} pairs with a stable
logical identifier (\texttt{functionRef}) resolved at deployment time. This is the indirection
layer that makes a workflow survive a function redeploy: the workflow references
\texttt{"fraud-check"}, and the concrete endpoint is bound later from the registry.

\paragraph{From a two-key registry to a metadata record.}
The registry model itself evolved. An early design keyed each function by two synthetic entries,
\texttt{"<functionRef>.host"} and \texttt{"<functionRef>.path"}, and resolved them through
\texttt{WorkflowInternalCallEndpointResolver} with the help of \texttt{FunctionEndpointKeys}.
This was superseded by a single structured record per function,
\texttt{FunctionInvocationMetadata} (\texttt{serviceName} plus a \texttt{url} that is an HTTPS
endpoint on GCP and a Lambda ARN on AWS), consumed by the two provider resolvers that actually
ship. As noted above, the older layer was never deleted --- it is the visible sediment of the
design's evolution, still exercised by tests but bypassed at runtime.

\paragraph{From single-provider to symmetric multi-provider.}
QuickFaaS originally targeted GCP (and Azure). The AWS Lambda provider added in this work
(Part~A) was deliberately shaped to mirror the GCP path rather than special-case it: the same
strategy interface, a parallel resolver, and a registry that stores an ARN in the same
\texttt{url} field a GCP HTTPS endpoint would occupy. The renderer distinguishes the two only at
the last step, emitting a native \texttt{lambda:invoke} task resource for internal AWS calls and
an \texttt{apigateway:invoke} HTTP task for external ones.

\paragraph{From naive to measured.}
Finally, the approach evolved under measurement. The registry began as a re-read-per-call file
(P3, P6), the resulting cost --- growing in proportion to both the number of calls and the
registry size --- was quantified, and the single-read optimisation
(P7) was designed as a direct response --- the clearest instance of the work using its own
benchmark suite to drive a concrete refinement, even if that refinement has yet to be
propagated into the production resolvers.

\section{Case Study: A Retail Bank (Santander Portugal)}
\label{sec:case-study-bank}

To make the trade-offs concrete, this section situates the integration inside a hypothetical
deployment at a large retail bank such as Santander Portugal. The scenarios below are
\emph{illustrative}; they use no proprietary or internal information and are grounded only in
publicly known regulatory and architectural constraints that any EU bank faces --- notably the
Digital Operational Resilience Act (DORA), the revised Payment Services Directive (PSD2/open
banking), and the least-privilege and change-control expectations of a regulated environment.

\subsection{Scenario 1 --- Redeploying a fraud-check function without touching workflows}

A fraud-scoring function sits on the critical path of card authorisation and is redeployed
frequently as models are retrained. In the pre-unification world, each redeploy that changed the
function's endpoint forced a corresponding edit to every workflow that called it --- a
change-heavy, error-prone coupling on a regulated payment flow. The Function Registry directly
targets this: workflows reference \texttt{internalFunction("fraud-check")}, and the endpoint is
rebound from the registry at deployment time. The bank gains exactly the property it wants ---
the orchestration definition is decoupled from the function's physical location, so a model
redeploy is an operations event, not a workflow-code change.

The current gaps show up here too. Because the registry is an unauthenticated file, the binding
that decides which endpoint a payment flow invokes is only as trustworthy as the file's
provenance; a bank would want that binding signed and stored in a controlled backend (see
\S\ref{sec:future-work}). And because the real resolvers still re-read the registry per call,
resolving a large workflow at deploy time is slower than it needs to be --- tolerable, but not
what one would ship.

\subsection{Scenario 2 --- Multi-region and multi-cloud resilience under DORA}

DORA pushes banks toward demonstrable operational resilience and credible exit strategies from
any single cloud provider. The symmetric GCP/AWS design is well aligned with this: the same
workflow definition, expressed once in the OmniFlow DSL, renders to AWS Step Functions or GCP
Cloud Workflows, and internal functions resolve to a Lambda ARN or a Cloud Run URL through the
identical registry abstraction. In principle the bank can keep one logical definition of a
process and target two providers, which is precisely the portability posture a resilience
regulator looks for.

The honest caveat is that portability is at the \emph{definition} layer, not the \emph{running}
layer: nothing here performs live cross-cloud failover, and the resolution cascade's regional
disambiguation (exact match, then \texttt{"region/functionRef"} suffix match, with a fail-fast
error on ambiguity) is a deployment-time convenience, not a runtime routing mechanism. It lowers
the cost of \emph{re-targeting}; it does not by itself deliver active-active resilience.

\subsection{Scenario 3 --- A regulator-facing audit trail}

Every registry write refreshes an \texttt{updatedAt} timestamp, and every resolution emits an
explicit, human-readable diagnostic trail (registry hit, provider discovery, or QuickFaaS
deploy). This is a useful foundation for the ``who bound what endpoint, and when'' question an
auditor asks. But a single mutable timestamp on a flat file is not an audit trail: there is no
history, no actor identity, and no tamper evidence. Meeting an audit obligation would require
turning the registry's write path into an append-only, attributable log --- again a hardening
task rather than a redesign, which speaks well of the underlying structure.

\subsection{What the case study reveals}

Across the three scenarios the same pattern recurs: the \emph{shape} of the solution is right
for a bank --- logical indirection, provider symmetry, explicit diagnostics --- while the
\emph{implementation maturity} is that of a thesis prototype. The gaps that matter in this
setting (registry integrity, least-privilege IAM, an attributable write history, propagating the
performance fix) are all incremental hardening of an already-sound design, not evidence that the
design is wrong. That is a favourable position for the contribution: it validates the concept and
leaves a well-scoped path to production.

\section{Future Work}
\label{sec:future-work}

The following items are ordered roughly by cost-to-benefit, from the smallest changes with the
clearest payoff to the larger structural ones.

\begin{enumerate}
    \item \textbf{Propagate the single-read optimisation into the real resolvers.} Convert
    \texttt{AwsInternalFunctionResolver} and \texttt{WorkflowInternalFunctionResolver} to read
    the registry once per \texttt{resolve} and thread the snapshot through, as
    \texttt{WorkflowInternalCallEndpointResolver} already does. This closes the gap between the
    measured optimum --- a cost proportional to the sum of the number of calls and the registry
    size, rather than their product --- and the still-multiplicative cost these classes pay today
    (P10, P11), and removes the incentive to keep the superseded resolver around.
    \item \textbf{Remove the design residue.} Delete
    \texttt{WorkflowInternalCallEndpointResolver} and \texttt{FunctionEndpointKeys} once their
    optimisation has been folded into the live resolvers, and align the error-semantics
    documentation with the shipped \texttt{serviceName}/\texttt{url} model.
    \item \textbf{Harden the registry.} Give the registry a pluggable backend behind its current
    interface --- a database or a secret/parameter store rather than a flat file --- with
    cross-process locking, an integrity guarantee (signing or checksums) on each entry, and an
    append-only, attributable write history to serve audit requirements. An in-memory cache with
    explicit invalidation would also fix the write-path scaling seen in P13/P17.
    \item \textbf{Replace the subprocess boundary.} Invoke QuickFaaS in-process as a library
    rather than shelling out to a JAR, removing the two-build-system coupling, the hard-coded
    path, and the fragile parsing of subprocess error text.
    \item \textbf{Move to least-privilege, change-controlled IAM.} Replace the automatic
    creation and mutation of shared IAM roles with explicit, pre-provisioned roles and a
    dry-run/plan mode that reports the IAM changes it \emph{would} make, so that privilege
    changes flow through the organisation's review pipeline instead of happening as a deployment
    side effect. Support an assume-role credential chain for per-tenant isolation.
    \item \textbf{Introduce test seams for the cloud-calling code.} Extract the AWS/GCP SDK
    interactions behind narrow interfaces so the deploy/IAM/readiness logic can be unit-tested
    against fakes, bringing the integration edges up to the coverage the pure core already
    enjoys --- without violating the project's rule against real cloud calls in tests.
    \item \textbf{Extend the provider set.} With the extension point validated on two providers,
    adding Azure Functions (already a QuickFaaS target) would exercise the strategy interface a
    third time and further substantiate the portability claim that matters for DORA-style exit
    strategies.
\end{enumerate}
